\documentclass[11pt]{scrartcl}

\usepackage[stixtwo]{fontsetup}
\usepackage{amsmath}
\usepackage{xfrac}

\begin{document}
\begin{enumerate}
	\item \textbf{Szöveges megoldás}

	      A számlálóban a leggyorsabban növekvő tag a $3 \cdot 2^{2n + 1}
		      = 6 \cdot 4^n$, hiszen
	      \[
		      n \ln n \ll n^2 \ll 4^n
		      \text.
	      \]
	      A nevezőben a leggyorsabban növekvő tag a $4^{n + 3} = 64 \cdot 4^n$,
	      hiszen
	      \[
		      n^6 \ll 4^n
		      \text.
	      \]
	      Mivel mind a számlálóban, mind a nevezőben a $c \cdot 4^n$ a domináns
	      tag, ezért a határérték ezen tagok együtthatóinak hányadosa lesz:
	      \[
		      \lim_{n \to \infty} \frac{
			      n \, \ln(n) + 3 \cdot 2^{2n + 1}
		      }{
			      4^{n + 3} + 5 n^6
		      } = \frac{6}{64} = \frac{3}{32}
		      \text.
	      \]

	      \textbf{Másik megoldás}

	      $3 \cdot 2^{2n + 1} = 6 \cdot 4^n$ és $4^{n + 3} = 64 \cdot 4^n$, így
	      \[
		      \frac{n \ln(n) + 6 \cdot 4^n}{64 \cdot 4^n + 5n^6}
		      = \frac{\frac{n \ln(n)}{4^n} + 6}{64 + \frac{5n^6}{4^n}}
		      \to \frac{6}{64}
		      = \frac{3}{32}
		      \text.
	      \]

	\item Behelyettesítéssel $(\infty - \infty)$ típusú határozatlan alakot
	      kapunk. Alkalmazzuk a gyökös kifejezések racionalizálását, vagyis
	      bővítsük a kifejezést a konjugáltjával:
	      \[
		      \sqrt{n^2 + 5n + 1} - n
		      = \frac{
			      (n^2 + 5n + 1) - n^2
		      }{
			      \sqrt{n^2 + 5n + 1} + n
		      }
		      = \frac{5n + 1}{\sqrt{n^2 + 5n + 1} + n}
		      \text.
	      \]
	      A határérték tehát:
	      \[
		      = \frac{5n + 1}{\sqrt{n^2 + 5n + 1} + n}
		      = \frac{5 + \sfrac{1}{n}}{\sqrt{1 + \sfrac{5}{n} + \sfrac{1}{n^2}} + 1}
		      \to \frac{5}{1 + 1} = \frac{5}{2}
		      \text.
	      \]

	\item \textbf{Számláló és nevező együtt kezelésével}
	      \[
		      \left( \frac{2n + 3}{5 + 2n} \right)^{n + 3}
		      = \left( \frac{2n + 5 - 2}{5 + 2n} \right)^{n + 3}
		      = \left[
			      \left( 1 + \frac{-2}{5 + 2n} \right)^{5 + 2n}
			      \right]^{\frac{n + 3}{5 + 2n}}
		      \to \left[e^{-2}\right]^{\frac{1}{2}} = e^{-1}
	      \]

	      \textbf{Számláló és nevező külön kezelésével}
	      \[
		      \left( \frac{2n + 3}{5 + 2n} \right)^{n + 3}
		      = \frac{\left(1 + \frac{3/2}{n}\right)^n}{\left(1 + \frac{5/2}{n}\right)^n}
		      \cdot \left( \frac{2n + 3}{2n + 5} \right)^3
		      \to \frac{e^{\sfrac{3}{2}}}{e^{\sfrac{5}{2}}} \cdot 1
		      = e^{-1}
	      \]
\end{enumerate}
\end{document}
